Restricting the test to the common support changes what the test measures: the comparison targets the region source and target share, not the full observed populations. Harmful change concentrated in weak-overlap regions will be downweighted and may go undetected. The procedure should not be a default replacement for unweighted testing; analysts should choose based on whether the monitoring question concerns the observed populations or only the comparable subpopulation.

Weight concentration is a finite-sample concern. Even with relative importance weighting, weak overlap can shrink effective sample size enough to make the common-support comparison noisy. ESS and power trade off continuously; lower $\lambda$ retains specificity but concentrates weights and reduces power. Practitioners should inspect the ESS diagnostic and the $\lambda$ sensitivity curve to balance false-alarm protection against power loss.

The method depends on domain-classifier quality. A miscalibrated, unstable, or poorly featured classifier can distort testing rather than improve it. Cross-validated or held-out domain probabilities help mitigate overfitting, and calibration checks on the domain classifier provide an additional quality signal. Appendix~\ref{app:domain-clf-sensitivity} shows that the doubly weighted test remains calibrated across logistic regression, random forest, and gradient-boosted classifiers in the tested settings, but this robustness should be validated per application.

This analysis covers non-sequential harmful-shift testing built from rank-consistent severity scores. The same common-support correction may apply to sequential risk trackers, canary tests, or process-control alarms, but those workflows introduce stopping-time and time-uniform error-control questions beyond our current scope.
